\documentclass[10pt,a4paper]{amsart}
\usepackage[margin=19mm]{geometry}
\usepackage{amsmath,amssymb,mathtools}
\usepackage[scr=rsfs,cal=euler]{mathalfa}
\usepackage{xfrac}
\usepackage[unicode,hidelinks,bookmarks=false]{hyperref}
\newcommand{\ie}{i.e.\ }
\newcommand{\cf}{cf.\ }
\newcommand{\resp}{respectively\ }
\newcommand{\Subsection}{\subsection}
\providecommand{\nobreakdash}{\nobreak\hbox{-}}
\setlength{\emergencystretch}{2em}
\multlinegap0pt
\usepackage{amssymb}
\usepackage{mathtools}
\usepackage[shortlabels]{enumitem}
\usepackage{tikz}
\usepackage{algorithm}\usepackage{algpseudocode}
\newtheorem{lemma}{Lemma}
\title{Betti numbers for cochordal zero-divisor graphs of commutative rings}
\author{Bilal Ahmad Rather}
\date{Source: arXiv:2605.13622v1 (CC BY 4.0)}
\begin{document}
\maketitle

\noindent\emph{Source and licence.} Betti numbers for cochordal zero-divisor graphs of commutative rings. Version arXiv:2605.13622v1 (CC BY 4.0), distributed by arXiv under the Creative Commons Attribution 4.0 International licence, http://creativecommons.org/licenses/by/4.0/, as recorded in the arXiv licence metadata for this exact version and shown on the abstract page https://arxiv.org/abs/2605.13622v1. Full licence notice: \texttt{licenses/arxiv-2605.13622-CC-BY-4.0.txt}.\par\smallskip

\noindent\textit{Editorial scope.} Counted unit: the article's lemma lem:valuation-annihilator (source0.tex lines 624--634). The article's complete original proof of it is delivered with it (source0.tex lines 637--702, one \texttt{proof} environment of 66 lines). No mathematics is written, repaired or re-derived: the statement and the proof are the article's own lines, in the article's own notation, and no retained line is substituted or re-flowed.  No further source-local context is needed: every cross-reference of every retained line resolves inside the two delivered spans.  Nothing was omitted from any retained span, so the package carries no omission record.  No retained line contains a citation, so the wrapper carries no bibliography and no bibliography entry.  No retained line contains a figure environment, an image-inclusion command or drawing code, so no image is delivered and the package declares no asset dependency.  Editorial formatting only, in addition: an AMS \texttt{amsart} preamble that restates the article's own declarations and macros used by the retained lines, namely the environments \texttt{lemma} on separate counters, together with the standard packages those lines need.  The retained environments print in this document as Lemma 1 in order of appearance, whereas the article prints the same items with the numbering of its own full document.  The complete native source of the article as distributed in the arXiv e-print for this exact version is bundled unchanged as \texttt{sources/200521/source0.tex}.
\begin{lemma}\label{lem:valuation-annihilator}
	Let $x,y\in \mathbb{Z}/n\mathbb{Z}$ be nonzero zero divisors.  Then $\operatorname{ann}(x)=\operatorname{ann}(y)$ if and only if $\nu_t(x)=\nu_t(y)$ for every $t$.  Moreover, the number of residues in the valuation class $\mathbf e=(e_1,\ldots,e_s)$ is $|C(\mathbf e)|=\prod_{t=1}^s c_t(e_t),$ where
	$$
	 c_t(e_t)=
	\begin{cases}
		p_t^{\alpha_t-e_t-1}(p_t-1),&0\leq e_t<\alpha_t,\\
		1,&e_t=\alpha_t.
	\end{cases}
	$$
	The all-maximal vector $(\alpha_1,\ldots,\alpha_s)$ corresponds to the zero residue and is omitted from the vertex set.
\end{lemma}

\begin{proof}
	Choose integer representatives of $x$ and $y$, with 
	$ n=\prod_{t=1}^{s}p_t^{\alpha_t}. $
	For a residue $x$, let
	$ d_x=\gcd(x,n). $
	Then
	$ d_x=\prod_{t=1}^{s}p_t^{\nu_t(x)}. $
	Clearly, the exponent of $p_t$ in $\gcd(x,n)$ is exactly the truncated
	$p_t$-adic valuation of $x$ modulo $n$. We first determine the annihilator of $x$.  A residue $z\in \mathbb{Z}/n\mathbb{Z}$ belongs
	to $\operatorname{ann}(x)$ if and only if
	 $xz\equiv 0 \pmod n, $
	or equivalently,
	$ n\mid xz,$  which is equivalent to
	$ \frac{n}{d_x}\mid z. $
	Hence
	$ \operatorname{ann}(x)=\left(\frac{n}{d_x}\right)
	\subseteq \mathbb{Z}/n\mathbb{Z}. $
	Therefore $\operatorname{ann}(x)$ is determined by $d_x$, and $d_x$ is determined by the
	valuation vector
	$ \nu(x)=(\nu_1(x),\ldots,\nu_s(x)). $
	Consequently, if $\nu_t(x)=\nu_t(y)$ for every $t$, then $d_x=d_y$, and hence
	$\operatorname{ann}(x)=\operatorname{ann}(y)$.
	
	 Conversely, suppose that $\operatorname{ann}(x)=\operatorname{ann}(y)$.  Since
	$ \operatorname{ann}(x)=\left(\frac{n}{d_x}\right)$ and $ \operatorname{ann}(y)=\left(\frac{n}{d_y}\right), $ so
	the two principal ideals generated by $n/d_x$ and $n/d_y$ in $\mathbb{Z}/n\mathbb{Z}$ are
	equal.  In $\mathbb{Z}/n\mathbb{Z}$, ideals are determined by divisors of $n$, so we have 
	$ \frac{n}{d_x}=\frac{n}{d_y}, $
	and hence $d_x=d_y$.  Comparing the prime-power decompositions of $d_x$ and
	$d_y$ gives
	$ \nu_t(x)=\nu_t(y) $ for every $t. $
	This proves the equivalence between equality of annihilators and equality of
	valuation vectors.
	
	 For counting the size of a valuation class, we use Chinese remainder theorem,
	$ \mathbb{Z}/n\mathbb{Z}\cong \prod_{t=1}^{s}\mathbb{Z}/p_t^{\alpha_t}\mathbb{Z}.$
	Thus the choices in different prime-power components are independent.  In the
	component $\mathbb{Z}/p_t^{\alpha_t}\mathbb{Z}$, the number of residues with exact
	valuation $e_t<\alpha_t$ is the number of multiples of $p_t^{e_t}$ which are
	not multiples of $p_t^{e_t+1}$, that is, such a number is
	$$
	p_t^{\alpha_t-e_t}-p_t^{\alpha_t-e_t-1}
	=
	p_t^{\alpha_t-e_t-1}(p_t-1).
	$$
	If $e_t=\alpha_t$, then there is only one residue with that valuation, namely
	the zero residue modulo $p_t^{\alpha_t}$. Therefore, for
	$ \mathbf e=(e_1,\ldots,e_s), $
	the number of residues with valuation vector $\mathbf e$ is
	$ |C(\mathbf e)|
	=
	\prod_{t=1}^{s}c_t(e_t), $
	where
	$$
	c_t(e_t)=
	\begin{cases}
		p_t^{\alpha_t-e_t-1}(p_t-1),&0\leq e_t<\alpha_t,\\
		1,&e_t=\alpha_t.
	\end{cases}
	$$
	Finally, the vector
	$ (\alpha_1,\ldots,\alpha_s) $
	means that every prime-power component is zero, and hence it represents the
	zero residue modulo $n$.  Since the zero-divisor graph uses only nonzero zero
	divisors as vertices, this class is omitted from the vertex set.
\end{proof}

\end{document}
